\scanpage{94}
\[
\begin{aligned}
\underline{H}^{11}(A\times_{S}B)&=\underline{H}^{11}(A/S)\oplus\underline{H}^{11}(B/S)\\
&\quad+\underline{H}^{1,0}(A/S)\otimes\underline{H}^{0,1}(B/S)+\underline{H}^{0,1}(A/S)\otimes\underline{H}^{1,0}(B/S)
\end{aligned}
\]
La classe $\unsure{\zeta|\xi}$ qu'on veut calculer a, pour des raisons de fonctorialité (regarder les inclusions $A\longrightarrow A\times_{S}B$, $B\longrightarrow A\times_{S}B$) des composantes nulles sur les deux premiers facteurs. Il reste à calculer sa classe dans le deuxième, on \unsure{trouve|devine|peut} ce qu'on \unsure{pense|fera|se faire}....

\begin{center}\rule{3cm}{0.35pt}\\[-0.6ex]\rule{4cm}{0.45pt}\end{center}

\deleted{Connexion \inserted{(absolue)} canonique sur $\underline{\mathcal M}_{X/S}$}

\deleted{\diagram{\begin{tikzcd}[ampersand replacement=\&,column sep=small]
X_{0}\arrow[r,dash]\arrow[d,dash] \& X'_{\unsure{1}}\arrow[d,dash] \& X''_{2}\arrow[dl,dash]\\
S_{0}\arrow[r,phantom,"\subset" description] \& S^{\unsure{\prime|' '}}
\end{tikzcd}}}
\margin{\deleted{dans le voisinage inf. du premier ordre de $S$ avec $\exists$ rétraction de $S$ sur $S_{0}$}}

\deleted{$\underline{W}_{X'_{1}/S'}\simeq\underline{W}_{X'_{2}/S'}$ canoniquement ??}

\deleted{\uline{Question équivalente} : Si on a un $L_{0}$ inversible sur $X_{0}$, muni d'une connexion\unsure{,} rel. à $S_{0}$, alors les prolongements (avec connexion) à $X'$ et à $X''$ se correspondent \unsure{biuniv.quement|biunivoquement}.}

$W_{\unsure{X/S}}$

$S_{0}\leftarrow S$

% archive sheet number: 93

